\thispagestyle{empty}

\begin{tikzpicture}[remember picture, overlay]
    
    % ===== 1. CABEÇALHO (topo) =====
    \node[anchor=north, yshift=-3cm, align=center] at (current page.north) {%
        \parbox{0.9\paperwidth}{%
            \centering
            \color{black}%
            \setlength{\parskip}{0pt}%
            \setlength{\baselineskip}{14pt}%
            {\fontspec{Times New Roman}\fontsize{12pt}{14pt}\bfseries\selectfont%
                \MakeUppercase{\universidade}}%
            \\[0pt]%
            {\fontspec{Times New Roman}\fontsize{12pt}{14pt}\selectfont%
                ``\universidadenome''}%
            \\[0pt]%
            {\fontspec{Times New Roman}\fontsize{12pt}{14pt}\selectfont%
                \institutenome}%
            \\[0pt]%
            {\fontspec{Times New Roman}\fontsize{12pt}{14pt}\selectfont%
                \campus}%
        }%
    };
    
    % ===== 2. NOME DO AUTOR =====
    \node[anchor=south, yshift=647pt, align=center] at (current page.south) {%
        \parbox{0.9\paperwidth}{%
            \centering
            \color{black}%
            {\fontspec{Times New Roman}\fontsize{12pt}{16pt}\bfseries\selectfont%
                \MakeUppercase{\autor}}%
        }%
    };
    
    % ===== 3. TÍTULO =====
    \node[anchor=south, yshift=590pt, align=center] at (current page.south) {%
        \parbox{0.9\paperwidth}{%
            \centering
            \color{black}%
            {\fontspec{Times New Roman}\fontsize{14pt}{16pt}\bfseries\selectfont%
                \MakeUppercase{\titulo}}%
        }%
    };
    
    % ===== 4. SUBTÍTULO =====
    \node[anchor=south, yshift=556pt, align=center] at (current page.south) {%
        \parbox{0.7\paperwidth}{%
            \centering
            \color{black}%
            {\fontspec{Times New Roman}\fontsize{14pt}{16pt}\bfseries\selectfont%
                \MakeUppercase{\subtitulo}}%
        }%
    };
    
    % ===== 5. TEXTO DE APRESENTAÇÃO =====
    \node[anchor=north, yshift=518pt, xshift=150pt, align=center] at (current page.south) {%
        \parbox{0.85\paperwidth}{%
            \centering
            \color{black}%
            \setlength{\parskip}{0pt}%
            \setlength{\baselineskip}{9pt}%
            \parbox{65mm}{%
                \justifying
                \hyphenpenalty=10000
                \exhyphenpenalty=10000
                \noindent
                {\fontspec{Times New Roman}\fontsize{12pt}{11pt}\selectfont%
                    \tipotxt\ de \graupor\ apresentada ao \institutenome\ %
                    do \campus, da \universidade\ ``\universidadenome'', %
                    como parte dos requisitos para obtenção do título de %
                    \grauautor\ em \area.%
                \par}%
            }%
        }%
    };
    
    % ===== 6. BANCA EXAMINADORA =====
    \node[anchor=north, yshift=380pt, align=center] at (current page.south) {%
        \parbox{0.85\paperwidth}{%
            \centering
            \color{black}%
            \setlength{\parskip}{0pt}%
            \setlength{\baselineskip}{12pt}%
            
            % Título da banca (centralizado)
            {\fontspec{Times New Roman}\fontsize{12pt}{14pt}\bfseries\selectfont%
                Comissão Examinadora
            \par}%
            \vspace{10pt}%
            
            % ----- MEMBRO 1: ORIENTADOR (PRESIDENTE) -----
            {\fontspec{Times New Roman}\fontsize{12pt}{14pt}\selectfont%
                \orientadortitulo\ \MakeUppercase{\orientador}%
            \par}%
            {\fontspec{Times New Roman}\fontsize{12pt}{12pt}\selectfont%
                \institutoabreviado / \universidadeabreviada / \cidade\ (\estado)%
            \par}%
            \vspace{10pt}%
            
            % ----- MEMBRO 2 -----
            {\fontspec{Times New Roman}\fontsize{12pt}{14pt}\selectfont%
                \titulomembrodois\ \MakeUppercase{\membrodois}%
            \par}%
            {\fontspec{Times New Roman}\fontsize{12pt}{12pt}\selectfont%
                \instmembrodois\ (\siglamembrodois)%
            \par}%
            \vspace{9pt}%
            
            % ----- MEMBRO 3 -----
            {\fontspec{Times New Roman}\fontsize{12pt}{14pt}\selectfont%
                \titulomembrotres\ \MakeUppercase{\membrotres}%
            \par}%
            {\fontspec{Times New Roman}\fontsize{12pt}{12pt}\selectfont%
                \instmembrotres\ (\siglamembrotres)%
            \par}%
            \vspace{9pt}%
            
            % ----- MEMBRO 4 -----
            {\fontspec{Times New Roman}\fontsize{12pt}{14pt}\selectfont%
                \titulomembroquatro\ \MakeUppercase{\membroquatro}%
            \par}%
            {\fontspec{Times New Roman}\fontsize{12pt}{12pt}\selectfont%
                \instmembroquatro\ (\siglamembroquatro)%
            \par}%
            \vspace{9pt}%
            
            % ----- MEMBRO 5 -----
            {\fontspec{Times New Roman}\fontsize{12pt}{14pt}\selectfont%
                \titulomembrocinco\ \MakeUppercase{\membrocinco}%
            \par}%
            {\fontspec{Times New Roman}\fontsize{12pt}{12pt}\selectfont%
                \instmembrocinco\ (\siglamembrocinco)%
            \par}%
        }%
    };
    
    % ===== 7. CONCEITO (centralizado com makebox) =====
    \node[anchor=south, yshift=114pt, align=center] at (current page.south) {%
        \makebox[\textwidth]{%
            \color{black}%
            \setlength{\parskip}{0pt}%
            \setlength{\baselineskip}{14pt}%
            {\fontspec{Times New Roman}\fontsize{12pt}{14pt}\selectfont%
                Conceito: \rule{3cm}{0.4pt}%
            }%
        }%
    };
    
    % ===== 8. CIDADE, ESTADO E DATA (rodapé) =====
    \node[anchor=south, yshift=70pt, align=center] at (current page.south) {%
        \parbox{0.9\paperwidth}{%
            \centering
            \color{black}%
            {\fontspec{Times New Roman}\fontsize{12pt}{14pt}\selectfont%
                \cidade\ (\estado), \datadefesaextenso%
            }%
        }%
    };
    
\end{tikzpicture}

\clearpage